% from ega1-6-fr.tex
\oldpage[I]{154}
on peut se borner au cas où $S=Y$ (\hyperref[I.3.3.11-fr]{3.3.11})~; posons
$f'=f_{(S')}$, et soit $U'$ un ouvert quasi-compact dans $S'$. Pour tout
$s'\in U'$, soit $T$ un voisinage ouvert affine de $g(s')$ dans $S$, et soit
$W$ un voisinage ouvert affine de $s'$ contenu dans
$U'\cap g^{-1}(T)$~; il suffira de montrer que $f'^{-1}(W)$ est
quasi-compact~; autrement dit, on peut se ramener à prouver que lorsque $S$
et $S'$ sont affines, l'espace sous-jacent à $X\times_S S'$ est
quasi-compact. Mais comme $X$ est alors par hypothèse réunion finie d'ouverts
affines $V_j$, $X\times_S S'$ est réunion des espaces sous-jacents aux
schémas affines $V_j\times_S S'$
(\hyperref[I.3.2.2-fr]{3.2.2} et \hyperref[I.3.2.7-fr]{3.2.7}), ce qui achève
de prouver la proposition.

On notera aussi que si $X=X'\amalg X''$ est somme de deux préschémas, un
morphisme $f:X\to Y$ est quasi-compact si et seulement si ses restrictions à
$X'$ et $X''$ le sont.

\begin{proposition}[6.6.5]
\label{I.6.6.5-fr}
Soit $f:X\to Y$ un morphisme quasi-compact. Pour que $f$ soit dominant, il
faut et il suffit que pour tout point générique $y$ d'une composante
irréductible de $Y$, $f^{-1}(y)$ contienne le point générique d'une composante
irréductible de $X$.
\end{proposition}

Il est immédiat que la condition est suffisante (sans supposer $f$
quasi-compact). Pour voir qu'elle est nécessaire, considérons un voisinage
ouvert affine $U$ de $y$~; $f^{-1}(U)$ est quasi-compact, donc réunion
\emph{finie} d'ouverts affines $V_i$, et l'hypothèse que $f$ est dominant
implique que $y$ appartient à l'adhérence dans $U$ d'un des $f(V_i)$. On peut
évidemment supposer $X$ et $Y$ réduits~; comme l'adhérence dans $X$ d'une
composante irréductible de $V_i$ est une composante irréductible de $X$
(\hyperref[0.2.1.6-fr]{0, 2.1.6}), on peut remplacer $X$ par $V_i$, $Y$ par
le sous-préschéma fermé réduit de $U$ ayant $\overline{f(V_i)}\cap U$ pour
espace sous-jacent (\hyperref[I.5.2.1-fr]{5.2.1}), et on est ainsi ramené à
démontrer la proposition lorsque $X=\operatorname{Spec}(A)$ et
$Y=\operatorname{Spec}(B)$ sont affines et réduits. Comme $f$ est dominant,
$B$ est alors un sous-anneau de $A$ (\hyperref[I.1.2.7-fr]{1.2.7}), et la
proposition résulte alors du fait que tout idéal premier minimal de $B$ est
l'intersection de $B$ et d'un idéal premier minimal de $A$
(\hyperref[0.1.5.8-fr]{0, 1.5.8}).

\begin{proposition}[6.6.6]
\label{I.6.6.6-fr}
\begin{enumerate}
  \item[(i)] Toute immersion locale est localement de type fini.
  \item[(ii)] Si deux morphismes $f:X\to Y$, $g:Y\to Z$ sont localement de
    type fini, il en est de même de $g\circ f$.
  \item[(iii)] Si $f:X\to Y$ est un $S$-morphisme localement de type fini,
    $f_{(S')}:X_{(S')}\to Y_{(S')}$ est localement de type fini pour toute
    extension $S'\to S$ du préschéma de base.
  \item[(iv)] Si $f:X\to X'$ et $g:Y\to Y'$ sont deux $S$-morphismes
    localement de type fini, $f\times_S g$ est localement de type fini.
  \item[(v)] Si le composé $g\circ f$ de deux morphismes est localement de
    type fini, $f$ est localement de type fini.
  \item[(vi)] Si un morphisme $f$ est localement de type fini, il en est de
    même de $f_{\mathrm{red}}$.
\end{enumerate}
\end{proposition}

En vertu de (\hyperref[I.5.5.12-fr]{5.5.12}), il suffit de démontrer (i), (ii)
et (iii). Si $j:X\to Y$ est une immersion locale, pour tout $x\in X$, il y a
un voisinage ouvert $V$ de $j(x)$ dans $Y$ et un voisinage ouvert $U$ de $x$
dans $X$ tels que la restriction de $j$ à $U$ soit une immersion fermée
$U\to V$ (\hyperref[I.4.5.1-fr]{4.5.1}), donc cette restriction est de type
fini. Pour établir (ii), considérons un point $x\in X$~; il y a par hypothèse
un voisinage ouvert $W$ de $g(f(x))$ et un voisinage ouvert $V$ de $f(x)$
tels que $g(V)\subset W$ et que $V$ soit de type fini sur $W$~; en outre
$f^{-1}(V)$ est localement de type fini sur $V$
(\hyperref[I.6.6.2-fr]{6.6.2}), donc il y a un voisinage ouvert $U$ de $x$ qui
